% !TEX root = ../tcc.tex

\begin{resumo}
O Ecossistema CEDIS, da Universidade de Brasília, reúne aplicações integradas de apoio ao ensino, mas suas análises permanecem restritas a indicadores descritivos de uso e contagens de respostas, sem estimar o domínio de conceitos dos estudantes. Este trabalho tem como objetivo desenvolver o Cognita, um módulo integrado ao ecossistema que aplica \textit{Bayesian Knowledge Tracing} (BKT) sobre respostas de múltiplas fontes para sustentar um ciclo duplo de \textit{feedback}: indicar aos estudantes os conceitos prioritários para revisão e sinalizar aos docentes os conceitos com baixa absorção na turma. A pesquisa, de natureza aplicada, abordagem quali-quantitativa e caráter exploratório, divide-se em duas fases. Na Fase~1 (TCC~1), executaram-se revisão narrativa da literatura, análise técnica do código-fonte do CEDIS Deck Manager, especificação e priorização de requisitos com \textit{Lean Inception} adaptada e modelo URPS+, projeto da arquitetura de integração, modelagem lógica relacional em conformidade com a Lei Geral de Proteção de Dados (LGPD) e prototipação de alta fidelidade. Como resultados desta etapa, obtiveram-se a caracterização dos dados e contratos do ecossistema, 25 requisitos funcionais, dos quais 20 compõem o Produto Mínimo Viável (MVP), e 17 requisitos não funcionais, a arquitetura e o modelo de dados com consentimento explícito e expurgo em cascata, dez protótipos validados e o cronograma do projeto. Conclui-se que o ecossistema reúne as condições técnicas para viabilizar o rastreamento de conhecimento interpretável e aderente à LGPD, consolidando a base para a Fase~2 (TCC~2), que abrangerá a submissão ao Comitê de Ética em Pesquisa, implementação do \textit{pipeline} BKT, avaliação preditiva comparada a um modelo de referência, validação docente e avaliação de usabilidade.

\vspace{\onelineskip}

\noindent
\textbf{Palavras-chave}: Rastreamento de Conhecimento. \textit{Bayesian Knowledge Tracing}. \textit{Learning Analytics}. Ecossistema CEDIS. Engenharia de Software.
\end{resumo}
